% chapters/conclusion.tex --- おわりに(番号あり)
\chapter{おわりに}

本研究では，HARの推論不確実度に基づくBLEアドバタイジング間隔の適応制御を対象とし，ルールベース方策の定義と評価指標を確立するとともに，提案手法の実機検証を行った．送信・受信・電力計測を統合した計測基盤を整備することで，$\si{\micro\coulomb}/\text{event}$と受信遅延に基づく期限超過率を同時に評価できる形に整理した．この結果，単純な閾値ルールでも固定間隔より省電力となる運用点が存在することを示し，将来のSafe Contextual Bandit最適化へ発展させるための基盤を構築した．

\section*{本研究のまとめ}
本研究で得られた知見を以下にまとめる．
\begin{itemize}
  \item 不確実度と安定度から複合スコアCCSを構成し，閾値・ヒステリシス・最小滞在時間を用いたルールベース方策として定義した．
  \item 送信TX，電力TXSD，受信RXの三ノード構成により，同一試行で電力とQoSを評価できる基盤を整備した．QoSはTLと$P_{\mathrm{out}}(\tau)$と定義する．
  \item TL/Poutの算出では，開始時刻のずれが結果を歪め得ることを示し，定数オフセット補正による時間同期手順として指標定義を固定した．
  \item オフライン評価により，固定間隔点と候補方策点のトレードオフを可視化し，実機実験の探索空間を縮小する枠組みを示した．
  \item Phase 1実測に基づくシミュレーションでGate B/Cを設け，環境シフト・非定常で破綻しにくい初期設定は$w=0$, reset=1, $m=0.01$とした．安全寄りの感度分析として$m=0.02/0.03$を併記した．
  \item 実機実験D2b/D3とアブレーションD4/D4Bにより，方策が固定間隔の間の運用点として成立することを確認し，さらに$U$が電力である短間隔滞在量を，CCSが同電力でQoSを改善する役割分担が存在することを示唆した．
  \item 高干渉環境E2では，CCSは固定\SI{100}{\milli\second}より低電力だが$P_{\\mathrm{out}}(2\\,\\mathrm{s})$/TLが悪化した．TailGuard追加により$P_{\\mathrm{out}}(2\\,\\mathrm{s})$を1.28\%，$TL_{p95}$を944.5 msまで改善し，QoS制約を満たす運用点を確認した．n=6とした．
\end{itemize}

\section*{今後の課題}
今後は，動的制御ではTL分布と$P_{\mathrm{out}}(\tau)$を高精度に算出する．これにより，オフライン評価の予測と実測の差分を体系的に説明できるよう整理する．また，行動集合を拡張し，制約$\epsilon$の境界に沿ってより良い運用点を探索する必要がある．さらに，Safe Contextual Banditによるオンライン最適化へ拡張し，環境変化でもQoS制約を維持しながら省電力化できる枠組みを検証する．









